\documentclass[12pt]{article}
\usepackage[utf8]{inputenc}
\usepackage[T1]{fontenc}
\usepackage[margin=1in]{geometry}
\usepackage{setspace}
\usepackage{xcolor}
\usepackage[colorlinks=true, urlcolor=blue!60!black, linkcolor=blue!60!black]{hyperref}
\setlength{\parindent}{2em}
\setlength{\parskip}{1.25em}
\renewcommand{\baselinestretch}{1.0}
\begin{document}
\begin{center}
{\LARGE \textit{Ashcroft v.\ Free Speech Coalition}, 535 U.S.\ 234 (2002)\par}
\vspace{0.8em}
{\large Case Brief --- Four-Page Edition\par}
\end{center}
\vspace{1em}

\section{Memory Jogger}

The Court struck down the Child Pornography Prevention Act of 1996's bans on sexually explicit images that merely ``appear to be'' of minors or ``convey the impression'' of depicting minors as substantially overbroad, holding that government may not suppress non-obscene speech created without harming real children (\href{https://supreme.justia.com/cases/federal/us/535/234/}{535 U.S.\ 234}).

\section{Detailed Case Facts}

Before 1996, federal law tracked \textit{New York v.\ Ferber} (1982), covering only sexually explicit depictions produced using actual minors. Concerned by advances in computer imaging, Congress enacted the CPPA, which kept that ban and added new categories carrying penalties of up to 15 years for a first non-possession offense. Section 2256(8)(B) reached any image that ``is, or appears to be, of a minor engaging in sexually explicit conduct,'' sweeping in wholly computer-generated images and youthful-looking adults without requiring obscenity under \textit{Miller v.\ California}. Section 2256(8)(D) reached any sexually explicit depiction marketed in a way that ``conveys the impression'' it depicts a minor, and the material stayed contraband in later possessors' hands. A ``morphing'' provision, \S~2256(8)(C), was not challenged. An affirmative defense for adult-only productions did not apply to possession or to computer-generated images.

Congress justified the ban on four grounds: virtual images could be used to seduce children, could whet pedophiles' appetites, could fuel the market for real child pornography, and could make it hard to prove that an image shows a real child. An adult-entertainment trade association, a nudist-lifestyle publisher, a painter, and a photographer brought a pre-enforcement facial challenge.

\section{Procedural History}

The district court granted the Government summary judgment, finding prosecution of legitimate works ``highly unlikely.'' The Ninth Circuit reversed, holding both provisions overbroad and vague. Because four other circuits had upheld the CPPA, the Supreme Court granted certiorari and affirmed on April 16, 2002.

\section{Judicial Votes}

Affirmed, \textbf{6--3} on virtual images under \S~2256(8)(B) and \textbf{7--2} on youthful-adult images and the pandering provision.
\begin{itemize}
\item \textbf{Majority:} Justice Kennedy, joined by Justices Stevens, Souter, Ginsburg, and Breyer.
\item \textbf{Concurring in the judgment:} Justice Thomas.
\item \textbf{Concurring in the judgment in part and dissenting in part:} Justice O'Connor; Chief Justice Rehnquist and Justice Scalia joined Part II.
\item \textbf{Dissenting:} Chief Justice Rehnquist, joined by Justice Scalia except as to one paragraph on legislative history.
\end{itemize}

\section{Holding}

Sections 2256(8)(B) and 2256(8)(D) are facially overbroad. Speech that is neither obscene under \textit{Miller} nor produced through the sexual exploitation of actual children under \textit{Ferber} retains First Amendment protection, and Congress may not ban it on the theory that it might indirectly encourage abuse or complicate prosecutions. A pandering statute is overbroad if it criminalizes possession of non-obscene adult works because someone earlier in the chain marketed them as child pornography. The Court did not reach vagueness; its remarks about morphed images and affirmative defenses are dicta that shaped the PROTECT Act of 2003.

\section{Analysis of Opinions}

\textbf{Majority (Kennedy).} Severe criminal penalties chill speech, making overbreadth review essential. The CPPA ignored \textit{Miller}'s three prongs, banning any single explicit scene of someone appearing under 18 and so threatening Renaissance paintings, \textit{Romeo and Juliet}, \textit{Traffic}, and \textit{American Beauty}. \textit{Ferber} and \textit{Osborne} allowed bans only because the speech was ``intrinsically related'' to real abuse; virtual images create no victims, and \textit{Ferber} itself pointed to youthful adults and simulation as lawful alternatives. The indirect-harm theories failed under \textit{Butler v.\ Michigan}, \textit{Stanley}, and \textit{Brandenburg}: the mere tendency of speech to encourage unlawful thoughts cannot justify suppression. The market theory was implausible, and banning protected speech because it resembles unprotected speech ``turns the First Amendment upside down.'' The affirmative defense shifted the burden to speakers and was unavailable for possession. The pandering provision, unlike \textit{Ginzburg}, required no commercial exploitation and tainted innocent downstream possessors.

\textbf{Thomas (concurring in the judgment).} The ``prosecution rationale'' was the Government's strongest argument, but it cited no acquittal caused by a computer-generated-image defense. If technology made real child pornography unprosecutable, Congress could enact a narrow restriction with an affirmative defense.

\textbf{O'Connor (in part).} She agreed that the pandering ban and the youthful-adult ban fail, but would uphold the ban on computer-generated images, reading ``appears to be'' as ``virtually indistinguishable from'' real children and deferring to Congress's predictive judgment.

\textbf{Rehnquist (dissenting).} He would save the whole statute by construction: ``sexually explicit conduct'' reaches only hard-core material, not actors ``squirming under a blanket,'' and the pandering clause reaches only the ``sordid business of pandering,'' with a strong \textit{mens rea} requirement for possession.

\section{Examples: Future Applications}

\begin{enumerate}
\item \textbf{Same-Side (Unconstitutional): Non-Obscene Anime Ban.} A State bans drawn or animated depictions of fictional characters who appears to be under 18'' in explicit conduct without \textit{Miller}'s prongs. No real child is abused, so the ban falls.
\item \textbf{Same-Side (Unconstitutional): Animal-Cruelty Depictions.} Relying on the case's refusal to create new unprotected categories, the Court struck down 18 U.S.C.\ \S~48 in \textit{United States v.\ Stevens} (2010).
\item \textbf{Opposite-Side (Constitutional): Offers of Actual Child Pornography.} The PROTECT Act's pandering provision punishes offers believed to involve an \textit{actual} minor, not downstream possession, and was upheld in \textit{United States v.\ Williams} (2008).
\item \textbf{Opposite-Side (Constitutional): Obscene Virtual Images.} A prosecution under 18 U.S.C.\ \S~1466A, which incorporates \textit{Miller}, is valid even if no real child was used.
\item \textbf{Fence-Sitter: Generative-AI Images.} Photorealistic AI images of nonexistent children depict no identifiable child, but they may realize Thomas's forensic deadlock, and a model trained on real abuse material arguably makes the output ``intrinsically related'' to real abuse.
\end{enumerate}

\section{Critique}

\textbf{Child-protection and technological critiques.} Scholars argued that the production-harm test understates real dangers, that the Court dismissed congressional findings on grooming and prosecution too cavalierly, and that its market assumption lacked empirical support. Generative AI has sharpened those critiques, as cheap photorealistic images threaten to overwhelm law enforcement.

\textbf{Doctrinal critiques.} Echoing Rehnquist and O'Connor, critics faulted the Court for using facial overbreadth rather than avoidance or severability and for relying on hypothetical prosecutions that never occurred, contrary to \textit{Broadrick}'s ``substantial'' overbreadth requirement. Joseph Blocher and others note that the case reinforced First Amendment categoricalism, praised for constraining balancing but strained by novel digital harms.

\section{Key Quotations}

\begin{quote}
``In contrast to the speech in \textit{Ferber}, speech that itself is the record of sexual abuse, the CPPA prohibits speech that records no crime and creates no victims by its production. Virtual child pornography is not `intrinsically related' to the sexual abuse of children, as were the materials in \textit{Ferber}.''
\end{quote}
\noindent --- Justice Kennedy, majority opinion, 535 U.S.\ at 250.

\begin{quote}
``The Government may not suppress lawful speech as the means to suppress unlawful speech. Protected speech does not become unprotected merely because it resembles the latter. The Constitution requires the reverse.''
\end{quote}
\noindent --- Justice Kennedy, majority opinion, 535 U.S.\ at 255.

\begin{quote}
``First Amendment freedoms are most in danger when the government seeks to control thought or to justify its laws for that impermissible end. The right to think is the beginning of freedom, and speech must be protected from the government because speech is the beginning of thought.''
\end{quote}
\noindent --- Justice Kennedy, majority opinion, 535 U.S.\ at 253.

\end{document}
